\section{Mechanisms and Alternative Explanations}
\label{sec:mechanisms}

We examine three mechanisms to distinguish the cover bidding hypothesis
from alternative explanations.

\subsection{M1: Strategic Selection vs.\ Competitive Displacement}

A key concern is whether FL presence \emph{causes} higher prices or
merely correlates with tender characteristics that independently drive
prices up. Under competitive displacement, FL firms would crowd out
genuine competitors, reducing effective competition. Under strategic
selection, cartels deploy cover bidders in tenders that already attract
genuine competition, and the cover bidders are added on top.

We regress the number of genuine (non-FL) competitors on FL presence.
The OLS coefficient is 0.143 ($p < 0.01$): FL-present tenders have
\emph{more}, not fewer, genuine competitors. This rules out competitive
displacement.

The IV estimate for non-FL firms (Table~\ref{tab:iv_main}, column~4) is
0.404 ($p < 0.01$), confirming the OLS finding: instrumented FL presence
is associated with substantially more genuine competitors. Cartels
appear to select into competitive tenders and add cover bidders on top
of existing competition.

\subsection{M2: Reference Price Calibration}

If cartels use cover bids to inflate reference prices for future
tenders, winners in FL-present tenders should bid closer to the
reference price. We regress $\log(p_{\text{negot}} / p_{\text{ref}})$ on
FL presence. The coefficient is $-0.041$ ($p < 0.01$), meaning
that winning bids in FL-present tenders are 4\% closer to the reference
price. This is consistent with cartels calibrating the winning bid near
the reference price while cover bidders submit bids that do not
undercut it.

\subsection{M3: Reverse Causality}

Could FL firms be attracted to high-priced tenders rather than causing
higher prices? We test this by regressing FL entry on \emph{lagged} log
prices within market (item $\times$ PBU) using a panel of 427,474
market-year observations. The coefficient is 0.0021 (SE~=~0.0008,
$p = 0.012$)---statistically significant but economically negligible.
In elasticity terms, a 10\% increase in lagged prices is associated with
a $0.1 \times 0.0021 = 0.00021$ increase in FL entry probability, or
0.02 percentage points against an unconditional FL rate of 4.8\%---an
elasticity of $0.00021 / 0.048 \approx 0.004$, essentially zero. This
small magnitude suggests that reverse causality---while detectable in a
sample of 1.6 million observations---cannot explain the 6--9\% price
effect attributed to FL presence.

\begin{table}[htbp]
\centering
\caption{Mechanism Tests}
\label{tab:mechanisms}
\begin{adjustbox}{max width=\textwidth}
\begin{threeparttable}
\small
\begin{tabular}{lccc}
\toprule
Mechanism & Coefficient & SE & N \\
\midrule
M1: Competitive displacement (OLS) & 0.1426*** & (0.0062) & 1,654,447 \\
\quad DV: $\log(n_{\text{genuine}} + 1)$ & & & \\
M1: Competitive displacement (IV) & 0.3244*** & (0.0249) & 1,654,447 \\
M2: Reference price calibration & -0.0412*** & (0.0031) & 1,654,394 \\
\quad DV: $\log(p_{\text{negot}} / p_{\text{ref}})$ & & & \\
M3: Reverse causality & 0.0021** & (0.0008) & 427,474 \\
\bottomrule
\end{tabular}
\begin{tablenotes}
\small
\item \textit{Notes:} M1: FL presence on genuine competition. M2: bid-to-reference ratio.
M3: lagged price predicting FL entry (reverse causality check).
*** \textit{p}$<$0.01, ** \textit{p}$<$0.05, * \textit{p}$<$0.1.
\end{tablenotes}
\end{threeparttable}
\end{adjustbox}
\end{table}


\subsection{Interpretation of the Convite--Preg\~ao Difference}

The FL coefficient is larger for preg\~ao (9.3\%) than for convite
(3.8\%). This may appear counterintuitive if preg\~ao's electronic
format provides more transparency. However, the difference reflects
two institutional features. First, preg\~ao allows real-time observation
of competing bids, which facilitates Regime~1 cover bidding: cover
bidders can monitor the auction and ensure their bids remain above the
designated winner's bid. Second, convite's sealed-bid format makes it
harder for cover bidders to calibrate bids---but also creates less need
for cover bidding when there are already three genuine bidders (the
minimum threshold). The smaller convite coefficient thus reflects both
reduced scope for cover bidding under sealed-bid format and the
compositional effect of the minimum-bidder rule.
